% Concept importance via insertion and deletion.
% Drop into Overleaf and % Concept importance via insertion and deletion.
% Drop into Overleaf and % Concept importance via insertion and deletion.
% Drop into Overleaf and % Concept importance via insertion and deletion.
% Drop into Overleaf and \input{concept_ablation} from the main document.
% Requires: booktabs, graphicx, xcolor (all standard).
% Bibliography entries are in concept_ablation.bib.

\section{Concept Importance via Insertion and Deletion}
\label{sec:concept-ablation}

TCAV~\cite{kim2018tcav} assigns each concept a sensitivity score, but a score is
not by itself evidence that the model \emph{uses} the concept. Insertion and
deletion~\cite{petsiuk2018rise} supplies that evidence: an explanation is
faithful if removing what it deems important degrades the prediction quickly,
and if restoring it to an uninformative baseline recovers the prediction
quickly. We adapt this protocol from input regions to concept directions.

\subsection{Method}
\label{sec:ablation-method}

Let $a \in \mathbb{R}^{D}$ be the activation of a sample at the layer in which
the Concept Activation Vectors (CAVs) were fitted, and let $c$ be the class the
model predicts for that sample. We orthonormalise the $K=12$ CAVs into a basis
$q_1,\dots,q_K$; orthonormalisation matters because the raw CAVs are strongly
correlated, and without it removing ``the top $k$ concepts'' would subtract a
shared component more than once.

For a ranking of the concepts, deletion and insertion at step $k$ are
\begin{align}
a^{\text{del}}_k &= a - \textstyle\sum_{i \in \mathcal{S}_k} (a \cdot q_i)\, q_i,\\
a^{\text{ins}}_k &= a - \textstyle\sum_{i=1}^{K} (a \cdot q_i)\, q_i
                     + \textstyle\sum_{i \in \mathcal{S}_k} (a \cdot q_i)\, q_i,
\end{align}
where $\mathcal{S}_k$ holds the $k$ highest-ranked concepts. By construction
$a^{\text{del}}_0 = a^{\text{ins}}_K = a$, so the two curves are complementary.
The modified activation is propagated through the remainder of the network and
we record the model's confidence in class $c$.

\paragraph{Confidence measure.}
We report the \emph{logit margin}, the logit of class $c$ minus the largest
competing logit, rather than the softmax probability. The classifier is
saturated on its own predictions (mean probability $0.99$), so probability-based
curves are flat to four decimal places regardless of the intervention; the
margin is a monotone transform of the same quantity but is not compressed.

\paragraph{Controls.}
Two controls separate two distinct claims. The \textsc{random} control applies
the identical procedure to the same concept subspace in shuffled order, testing
whether the \emph{ranking} carries information. The \textsc{noise} control
replaces the CAVs with $K$ random orthonormal directions, testing whether the
\emph{subspace} carries information or whether any $K$-dimensional subspace
would behave the same. All experiments use the $1{,}335$ correctly classified
RAVDESS recordings.

\subsection{Faithfulness of the TCAV ranking}
\label{sec:ablation-faithfulness}

Table~\ref{tab:ablation-auc} reports the area under the deletion and insertion
curves. A faithful ranking yields a low deletion AUC and a high insertion AUC.

\begin{table}[t]
\centering
\caption{Insertion--deletion areas under the curve, and the total confidence
lost when all $K=12$ directions are removed. Lower deletion AUC and higher
insertion AUC indicate a more faithful ranking.}
\label{tab:ablation-auc}
\begin{tabular}{lccc}
\toprule
Ordering & Deletion AUC $\downarrow$ & Insertion AUC $\uparrow$ & Total drop \\
\midrule
TCAV importance      & 8.82 & 8.67 & $-0.769$ \\
Shuffled (control)   & 8.75 & 8.75 & $-0.769$ \\
Random subspace (noise) & 9.09 & 9.09 & $-0.002$ \\
\bottomrule
\end{tabular}
\end{table}

Two conclusions follow, and they point in opposite directions.

The concepts are informative \emph{as a set}. Deleting the $12$ concept
directions costs $0.769$ of margin, whereas deleting $12$ random directions
costs $0.002$, a factor of roughly $380$. Since both interventions remove
exactly $12$ of $D=71{,}680$ dimensions, the difference is attributable
entirely to \emph{which} subspace is removed: the CAVs are aligned with
class-relevant structure rather than with arbitrary directions.

The TCAV ranking, however, is not faithful. Ordering deletions by TCAV
importance (AUC $8.82$) is no better than shuffling them (AUC $8.75$), and both
necessarily converge to the same endpoint because they remove the same subspace.
Were the per-sample importance scores meaningful, removing the highest-ranked
concepts first would degrade the prediction faster than removing them in
arbitrary order. It does not.

This is consistent with the geometry of the CAVs themselves: fitted pairwise
against a random counterexample set, the twelve directions have a mean pairwise
cosine similarity of $0.995$, so they span approximately one effective
dimension. A ranking cannot recover per-concept structure that the
representation does not contain.

\subsection{Magnitude of the effect}
\label{sec:ablation-percent}

The areas in Table~\ref{tab:ablation-auc} are in logit units, which are hard to
interpret on their own. Normalising every curve by the untouched decision margin
expresses the same result as a percentage of the model's confidence
(Figure~\ref{fig:percent-impact}).

Removing all twelve concept directions costs $8.5\%$ of the decision margin;
removing twelve random directions costs $0.02\%$. The concepts are therefore
roughly $380$ times more consequential than an arbitrary subspace of the same
dimension, while still accounting for under a tenth of the evidence the model
uses. This is the quantitative form of the qualified conclusion above: the
concept vocabulary is real but partial, and the remaining $91.5\%$ of the margin
rests on structure the twelve hand-designed patterns do not describe.

\begin{figure*}[t]
\centering
\includegraphics[width=\linewidth]{percent_impact.png}
\caption{Insertion and deletion expressed as a percentage of the original
decision margin. Left: deleting concept directions, most important first.
Centre: restoring them to a stripped activation. Right: the share of the margin
carried by the concepts, per emotion. Negative values indicate emotions whose
confidence \emph{increases} when the concepts are removed.}
\label{fig:percent-impact}
\end{figure*}

\begin{table}[t]
\centering
\caption{Share of the decision margin carried by the twelve concepts, per
emotion (ratio of means; $n$ is the number of correctly classified recordings).
A negative value means the concepts collectively argue \emph{against} the
predicted class.}
\label{tab:percent-emotion}
\begin{tabular}{lrrrr}
\toprule
Emotion & $n$ & Baseline & After removal & \% of margin \\
\midrule
sad       & 168 & 8.44 & 5.74 & $+32.0$ \\
disgusted & 182 & 9.09 & 6.72 & $+26.1$ \\
happy     & 169 & 9.14 & 8.54 & $+6.7$ \\
surprised & 180 & 8.67 & 8.09 & $+6.7$ \\
angry     & 182 & 8.57 & 8.01 & $+6.5$ \\
neutral   &  92 & 9.14 & 8.55 & $+6.4$ \\
calm      & 184 & 9.69 & 9.47 & $+2.3$ \\
fearful   & 178 & 9.94 & 11.41 & $-14.7$ \\
\midrule
overall   & 1335 & 9.09 & 8.32 & $+8.5$ \\
\bottomrule
\end{tabular}
\end{table}

The per-emotion split in Table~\ref{tab:percent-emotion} is markedly uneven. The
concepts carry a third of the margin for \emph{sad} and a quarter for
\emph{disgusted}, but only a few percent for \emph{calm}, \emph{angry},
\emph{neutral}, \emph{happy} and \emph{surprised}. For \emph{fearful} the effect
is negative: removing the concepts \emph{raises} the margin by $14.7\%$, so on
balance the modelled patterns are evidence against that class rather than for
it. The vocabulary is thus most explanatory for the sustained, low-arousal
classes and least explanatory---indeed counter-explanatory---for \emph{fearful}.

\subsection{Per-concept importance}
\label{sec:ablation-per-concept}

Because the ranking is uninformative, we attribute importance directly by
leave-one-out deletion: each concept is removed on its own and we record the
resulting change in margin. Table~\ref{tab:concept-importance} gives the result
aggregated over all emotions.

\begin{table}[t]
\centering
\caption{Leave-one-out concept importance, averaged over all correctly
classified recordings. A single random direction is included for scale.}
\label{tab:concept-importance}
\begin{tabular}{lcc}
\toprule
Concept & Margin drop & Share of total \\
\midrule
\texttt{long\_constant\_thick}        & 0.247 & 53.2\% \\
\texttt{long\_dropping\_steep\_thick} & 0.089 & 19.2\% \\
\texttt{short\_dropping\_steep\_thin} & 0.061 & 13.2\% \\
\texttt{short\_rising\_steep\_thick}  & 0.016 & 3.4\% \\
\texttt{long\_dropping\_flat\_thick}  & 0.014 & 3.0\% \\
\texttt{short\_rising\_steep\_thin}   & 0.010 & 2.1\% \\
\texttt{long\_dropping\_steep\_thin}  & 0.008 & 1.7\% \\
\texttt{short\_dropping\_steep\_thick}& 0.006 & 1.3\% \\
\texttt{long\_rising\_flat\_thick}    & 0.006 & 1.3\% \\
\texttt{long\_rising\_steep\_thin}    & 0.003 & 0.7\% \\
\texttt{long\_rising\_steep\_thick}   & 0.003 & 0.6\% \\
\texttt{short\_constant\_thick}       & 0.002 & 0.4\% \\
\midrule
\emph{single random direction}        & 0.0004 & --- \\
\bottomrule
\end{tabular}
\end{table}

Importance is heavily concentrated. \texttt{long\_constant\_thick}, a long
horizontal band corresponding to sustained harmonic energy, accounts for over
half of the total effect; the three highest-ranked concepts account for $85\%$.
The six lowest-ranked concepts contribute under $6\%$ combined, and the weakest
is within an order of magnitude of a random direction. We note that the
individual drops sum to $0.46$ while removing all twelve jointly costs $0.74$:
the concepts interact through the network rather than contributing additively,
so these values constitute a ranking rather than a decomposition.

\subsection{Per-emotion analysis}
\label{sec:ablation-per-emotion}

Aggregating over emotions conceals the more informative structure. Repeating the
leave-one-out analysis within each emotion (Table~\ref{tab:importance-emotion},
Figure~\ref{fig:importance-emotion}) shows that a concept's effect is
\emph{signed}: a positive drop means the concept supported the prediction, and a
negative drop means it argued against it, so that removing it increased the
model's confidence.

\begin{table}[t]
\centering
\small
\setlength{\tabcolsep}{3.5pt}
\caption{Change in logit margin when each concept is deleted, by emotion.
Positive values indicate the concept supports the prediction; negative values
indicate it opposes it. Concept names are abbreviated: L/S = long/short,
const/drop/rise = constant/dropping/rising, st/fl = steep/flat,
tk/tn = thick/thin.}
\label{tab:importance-emotion}
\begin{tabular}{lrrrrrrrrrrrr}
\toprule
& \rotatebox{90}{L-const/tk} & \rotatebox{90}{L-drop/fl/tk} & \rotatebox{90}{L-drop/st/tk}
& \rotatebox{90}{L-drop/st/tn} & \rotatebox{90}{L-rise/fl/tk} & \rotatebox{90}{L-rise/st/tk}
& \rotatebox{90}{L-rise/st/tn} & \rotatebox{90}{S-const/tk} & \rotatebox{90}{S-drop/st/tk}
& \rotatebox{90}{S-drop/st/tn} & \rotatebox{90}{S-rise/st/tk} & \rotatebox{90}{S-rise/st/tn} \\
\midrule
angry     &  0.27 & $-0.39$ &  0.33 &  0.04 & $-0.15$ & $-0.02$ &  0.02 & $-0.06$ & $-0.10$ &  0.25 &  0.09 & $-0.00$ \\
calm      & $-0.19$ &  0.05 & $-0.18$ &  0.02 &  0.20 & $-0.00$ & $-0.00$ & $-0.02$ &  0.08 & $-0.06$ &  0.01 &  0.04 \\
disgusted &  1.09 &  0.33 &  0.23 &  0.00 &  0.15 &  0.05 &  0.00 &  0.07 &  0.06 &  0.11 & $-0.05$ & $-0.02$ \\
fearful   & $-0.86$ & $-0.37$ & $-0.26$ & $-0.00$ & $-0.15$ & $-0.04$ & $-0.00$ & $-0.07$ & $-0.13$ &  0.14 &  0.06 &  0.02 \\
happy     &  0.49 & $-0.19$ &  0.01 & $-0.01$ & $-0.12$ & $-0.02$ &  0.01 &  0.01 & $-0.03$ &  0.05 &  0.04 &  0.01 \\
neutral   &  0.07 &  0.08 &  0.10 & $-0.00$ & $-0.11$ & $-0.03$ &  0.00 &  0.03 &  0.15 &  0.10 & $-0.04$ &  0.02 \\
sad       &  1.05 &  0.65 &  0.24 &  0.03 &  0.34 &  0.07 & $-0.00$ &  0.05 &  0.09 & $-0.22$ &  0.02 &  0.02 \\
surprised &  0.15 &  0.08 &  0.18 & $-0.00$ & $-0.01$ &  0.02 & $-0.00$ &  0.01 & $-0.03$ &  0.05 & $-0.02$ & $-0.02$ \\
\bottomrule
\end{tabular}
\end{table}

\begin{figure}[t]
\centering
\includegraphics[width=\linewidth]{importance_by_emotion.png}
\caption{Per-emotion concept importance. Red indicates that deleting the concept
reduces confidence (the concept supports the prediction); blue indicates that
deleting it increases confidence (the concept opposes the prediction).}
\label{fig:importance-emotion}
\end{figure}

The dominant concept, \texttt{long\_constant\_thick}, does not act uniformly. It
is the strongest positive contributor for \emph{sad} ($+1.05$),
\emph{disgusted} ($+1.09$) and \emph{happy} ($+0.49$), yet the strongest
negative contributor for \emph{fearful} ($-0.86$) and a negative contributor for
\emph{calm} ($-0.19$). Sustained harmonic energy therefore functions as
discriminative evidence in both directions: its presence pushes the model toward
the sustained-phonation emotions and away from \emph{fearful}, whose
realisations are characteristically less steady. An aggregate analysis averages
these opposing roles toward zero and loses the effect entirely.

Two further patterns are visible. For \emph{angry}, the falling-tone concepts
separate by steepness in opposite directions: \texttt{long\_dropping\_steep\_thick}
supports the prediction ($+0.33$) while \texttt{long\_dropping\_flat\_thick}
opposes it ($-0.39$), consistent with anger being marked by sharp rather than
gradual pitch falls. For \emph{fearful}, eleven of the twelve concepts have
negative or negligible effect, indicating that the class is identified largely
by the \emph{absence} of the modelled spectrogram structures rather than by the
presence of any of them.

\subsection{Discussion}
\label{sec:ablation-discussion}

The ablation study supports a qualified conclusion. The concept set is
meaningful: the subspace it spans is strongly class-relevant relative to a
random subspace of equal dimension, and leave-one-out deletion identifies a
clear and acoustically interpretable ordering, dominated by sustained harmonic
energy. The per-emotion decomposition further recovers signed, emotion-specific
roles that are consistent with the phonetic character of the classes.

The TCAV sensitivity scores themselves, however, do not survive the faithfulness
test: their ranking performs no better than chance under insertion and deletion.
We attribute this to the near-collinearity of the pairwise-fitted CAVs. This
suggests that concept-based explanation of this model should report ablation
importance rather than TCAV magnitude, and that the concept vocabulary would
benefit from being fitted so as to be mutually discriminative rather than merely
distinguishable from a random baseline.
 from the main document.
% Requires: booktabs, graphicx, xcolor (all standard).
% Bibliography entries are in concept_ablation.bib.

\section{Concept Importance via Insertion and Deletion}
\label{sec:concept-ablation}

TCAV~\cite{kim2018tcav} assigns each concept a sensitivity score, but a score is
not by itself evidence that the model \emph{uses} the concept. Insertion and
deletion~\cite{petsiuk2018rise} supplies that evidence: an explanation is
faithful if removing what it deems important degrades the prediction quickly,
and if restoring it to an uninformative baseline recovers the prediction
quickly. We adapt this protocol from input regions to concept directions.

\subsection{Method}
\label{sec:ablation-method}

Let $a \in \mathbb{R}^{D}$ be the activation of a sample at the layer in which
the Concept Activation Vectors (CAVs) were fitted, and let $c$ be the class the
model predicts for that sample. We orthonormalise the $K=12$ CAVs into a basis
$q_1,\dots,q_K$; orthonormalisation matters because the raw CAVs are strongly
correlated, and without it removing ``the top $k$ concepts'' would subtract a
shared component more than once.

For a ranking of the concepts, deletion and insertion at step $k$ are
\begin{align}
a^{\text{del}}_k &= a - \textstyle\sum_{i \in \mathcal{S}_k} (a \cdot q_i)\, q_i,\\
a^{\text{ins}}_k &= a - \textstyle\sum_{i=1}^{K} (a \cdot q_i)\, q_i
                     + \textstyle\sum_{i \in \mathcal{S}_k} (a \cdot q_i)\, q_i,
\end{align}
where $\mathcal{S}_k$ holds the $k$ highest-ranked concepts. By construction
$a^{\text{del}}_0 = a^{\text{ins}}_K = a$, so the two curves are complementary.
The modified activation is propagated through the remainder of the network and
we record the model's confidence in class $c$.

\paragraph{Confidence measure.}
We report the \emph{logit margin}, the logit of class $c$ minus the largest
competing logit, rather than the softmax probability. The classifier is
saturated on its own predictions (mean probability $0.99$), so probability-based
curves are flat to four decimal places regardless of the intervention; the
margin is a monotone transform of the same quantity but is not compressed.

\paragraph{Controls.}
Two controls separate two distinct claims. The \textsc{random} control applies
the identical procedure to the same concept subspace in shuffled order, testing
whether the \emph{ranking} carries information. The \textsc{noise} control
replaces the CAVs with $K$ random orthonormal directions, testing whether the
\emph{subspace} carries information or whether any $K$-dimensional subspace
would behave the same. All experiments use the $1{,}335$ correctly classified
RAVDESS recordings.

\subsection{Faithfulness of the TCAV ranking}
\label{sec:ablation-faithfulness}

Table~\ref{tab:ablation-auc} reports the area under the deletion and insertion
curves. A faithful ranking yields a low deletion AUC and a high insertion AUC.

\begin{table}[t]
\centering
\caption{Insertion--deletion areas under the curve, and the total confidence
lost when all $K=12$ directions are removed. Lower deletion AUC and higher
insertion AUC indicate a more faithful ranking.}
\label{tab:ablation-auc}
\begin{tabular}{lccc}
\toprule
Ordering & Deletion AUC $\downarrow$ & Insertion AUC $\uparrow$ & Total drop \\
\midrule
TCAV importance      & 8.82 & 8.67 & $-0.769$ \\
Shuffled (control)   & 8.75 & 8.75 & $-0.769$ \\
Random subspace (noise) & 9.09 & 9.09 & $-0.002$ \\
\bottomrule
\end{tabular}
\end{table}

Two conclusions follow, and they point in opposite directions.

The concepts are informative \emph{as a set}. Deleting the $12$ concept
directions costs $0.769$ of margin, whereas deleting $12$ random directions
costs $0.002$, a factor of roughly $380$. Since both interventions remove
exactly $12$ of $D=71{,}680$ dimensions, the difference is attributable
entirely to \emph{which} subspace is removed: the CAVs are aligned with
class-relevant structure rather than with arbitrary directions.

The TCAV ranking, however, is not faithful. Ordering deletions by TCAV
importance (AUC $8.82$) is no better than shuffling them (AUC $8.75$), and both
necessarily converge to the same endpoint because they remove the same subspace.
Were the per-sample importance scores meaningful, removing the highest-ranked
concepts first would degrade the prediction faster than removing them in
arbitrary order. It does not.

This is consistent with the geometry of the CAVs themselves: fitted pairwise
against a random counterexample set, the twelve directions have a mean pairwise
cosine similarity of $0.995$, so they span approximately one effective
dimension. A ranking cannot recover per-concept structure that the
representation does not contain.

\subsection{Magnitude of the effect}
\label{sec:ablation-percent}

The areas in Table~\ref{tab:ablation-auc} are in logit units, which are hard to
interpret on their own. Normalising every curve by the untouched decision margin
expresses the same result as a percentage of the model's confidence
(Figure~\ref{fig:percent-impact}).

Removing all twelve concept directions costs $8.5\%$ of the decision margin;
removing twelve random directions costs $0.02\%$. The concepts are therefore
roughly $380$ times more consequential than an arbitrary subspace of the same
dimension, while still accounting for under a tenth of the evidence the model
uses. This is the quantitative form of the qualified conclusion above: the
concept vocabulary is real but partial, and the remaining $91.5\%$ of the margin
rests on structure the twelve hand-designed patterns do not describe.

\begin{figure*}[t]
\centering
\includegraphics[width=\linewidth]{percent_impact.png}
\caption{Insertion and deletion expressed as a percentage of the original
decision margin. Left: deleting concept directions, most important first.
Centre: restoring them to a stripped activation. Right: the share of the margin
carried by the concepts, per emotion. Negative values indicate emotions whose
confidence \emph{increases} when the concepts are removed.}
\label{fig:percent-impact}
\end{figure*}

\begin{table}[t]
\centering
\caption{Share of the decision margin carried by the twelve concepts, per
emotion (ratio of means; $n$ is the number of correctly classified recordings).
A negative value means the concepts collectively argue \emph{against} the
predicted class.}
\label{tab:percent-emotion}
\begin{tabular}{lrrrr}
\toprule
Emotion & $n$ & Baseline & After removal & \% of margin \\
\midrule
sad       & 168 & 8.44 & 5.74 & $+32.0$ \\
disgusted & 182 & 9.09 & 6.72 & $+26.1$ \\
happy     & 169 & 9.14 & 8.54 & $+6.7$ \\
surprised & 180 & 8.67 & 8.09 & $+6.7$ \\
angry     & 182 & 8.57 & 8.01 & $+6.5$ \\
neutral   &  92 & 9.14 & 8.55 & $+6.4$ \\
calm      & 184 & 9.69 & 9.47 & $+2.3$ \\
fearful   & 178 & 9.94 & 11.41 & $-14.7$ \\
\midrule
overall   & 1335 & 9.09 & 8.32 & $+8.5$ \\
\bottomrule
\end{tabular}
\end{table}

The per-emotion split in Table~\ref{tab:percent-emotion} is markedly uneven. The
concepts carry a third of the margin for \emph{sad} and a quarter for
\emph{disgusted}, but only a few percent for \emph{calm}, \emph{angry},
\emph{neutral}, \emph{happy} and \emph{surprised}. For \emph{fearful} the effect
is negative: removing the concepts \emph{raises} the margin by $14.7\%$, so on
balance the modelled patterns are evidence against that class rather than for
it. The vocabulary is thus most explanatory for the sustained, low-arousal
classes and least explanatory---indeed counter-explanatory---for \emph{fearful}.

\subsection{Per-concept importance}
\label{sec:ablation-per-concept}

Because the ranking is uninformative, we attribute importance directly by
leave-one-out deletion: each concept is removed on its own and we record the
resulting change in margin. Table~\ref{tab:concept-importance} gives the result
aggregated over all emotions.

\begin{table}[t]
\centering
\caption{Leave-one-out concept importance, averaged over all correctly
classified recordings. A single random direction is included for scale.}
\label{tab:concept-importance}
\begin{tabular}{lcc}
\toprule
Concept & Margin drop & Share of total \\
\midrule
\texttt{long\_constant\_thick}        & 0.247 & 53.2\% \\
\texttt{long\_dropping\_steep\_thick} & 0.089 & 19.2\% \\
\texttt{short\_dropping\_steep\_thin} & 0.061 & 13.2\% \\
\texttt{short\_rising\_steep\_thick}  & 0.016 & 3.4\% \\
\texttt{long\_dropping\_flat\_thick}  & 0.014 & 3.0\% \\
\texttt{short\_rising\_steep\_thin}   & 0.010 & 2.1\% \\
\texttt{long\_dropping\_steep\_thin}  & 0.008 & 1.7\% \\
\texttt{short\_dropping\_steep\_thick}& 0.006 & 1.3\% \\
\texttt{long\_rising\_flat\_thick}    & 0.006 & 1.3\% \\
\texttt{long\_rising\_steep\_thin}    & 0.003 & 0.7\% \\
\texttt{long\_rising\_steep\_thick}   & 0.003 & 0.6\% \\
\texttt{short\_constant\_thick}       & 0.002 & 0.4\% \\
\midrule
\emph{single random direction}        & 0.0004 & --- \\
\bottomrule
\end{tabular}
\end{table}

Importance is heavily concentrated. \texttt{long\_constant\_thick}, a long
horizontal band corresponding to sustained harmonic energy, accounts for over
half of the total effect; the three highest-ranked concepts account for $85\%$.
The six lowest-ranked concepts contribute under $6\%$ combined, and the weakest
is within an order of magnitude of a random direction. We note that the
individual drops sum to $0.46$ while removing all twelve jointly costs $0.74$:
the concepts interact through the network rather than contributing additively,
so these values constitute a ranking rather than a decomposition.

\subsection{Per-emotion analysis}
\label{sec:ablation-per-emotion}

Aggregating over emotions conceals the more informative structure. Repeating the
leave-one-out analysis within each emotion (Table~\ref{tab:importance-emotion},
Figure~\ref{fig:importance-emotion}) shows that a concept's effect is
\emph{signed}: a positive drop means the concept supported the prediction, and a
negative drop means it argued against it, so that removing it increased the
model's confidence.

\begin{table}[t]
\centering
\small
\setlength{\tabcolsep}{3.5pt}
\caption{Change in logit margin when each concept is deleted, by emotion.
Positive values indicate the concept supports the prediction; negative values
indicate it opposes it. Concept names are abbreviated: L/S = long/short,
const/drop/rise = constant/dropping/rising, st/fl = steep/flat,
tk/tn = thick/thin.}
\label{tab:importance-emotion}
\begin{tabular}{lrrrrrrrrrrrr}
\toprule
& \rotatebox{90}{L-const/tk} & \rotatebox{90}{L-drop/fl/tk} & \rotatebox{90}{L-drop/st/tk}
& \rotatebox{90}{L-drop/st/tn} & \rotatebox{90}{L-rise/fl/tk} & \rotatebox{90}{L-rise/st/tk}
& \rotatebox{90}{L-rise/st/tn} & \rotatebox{90}{S-const/tk} & \rotatebox{90}{S-drop/st/tk}
& \rotatebox{90}{S-drop/st/tn} & \rotatebox{90}{S-rise/st/tk} & \rotatebox{90}{S-rise/st/tn} \\
\midrule
angry     &  0.27 & $-0.39$ &  0.33 &  0.04 & $-0.15$ & $-0.02$ &  0.02 & $-0.06$ & $-0.10$ &  0.25 &  0.09 & $-0.00$ \\
calm      & $-0.19$ &  0.05 & $-0.18$ &  0.02 &  0.20 & $-0.00$ & $-0.00$ & $-0.02$ &  0.08 & $-0.06$ &  0.01 &  0.04 \\
disgusted &  1.09 &  0.33 &  0.23 &  0.00 &  0.15 &  0.05 &  0.00 &  0.07 &  0.06 &  0.11 & $-0.05$ & $-0.02$ \\
fearful   & $-0.86$ & $-0.37$ & $-0.26$ & $-0.00$ & $-0.15$ & $-0.04$ & $-0.00$ & $-0.07$ & $-0.13$ &  0.14 &  0.06 &  0.02 \\
happy     &  0.49 & $-0.19$ &  0.01 & $-0.01$ & $-0.12$ & $-0.02$ &  0.01 &  0.01 & $-0.03$ &  0.05 &  0.04 &  0.01 \\
neutral   &  0.07 &  0.08 &  0.10 & $-0.00$ & $-0.11$ & $-0.03$ &  0.00 &  0.03 &  0.15 &  0.10 & $-0.04$ &  0.02 \\
sad       &  1.05 &  0.65 &  0.24 &  0.03 &  0.34 &  0.07 & $-0.00$ &  0.05 &  0.09 & $-0.22$ &  0.02 &  0.02 \\
surprised &  0.15 &  0.08 &  0.18 & $-0.00$ & $-0.01$ &  0.02 & $-0.00$ &  0.01 & $-0.03$ &  0.05 & $-0.02$ & $-0.02$ \\
\bottomrule
\end{tabular}
\end{table}

\begin{figure}[t]
\centering
\includegraphics[width=\linewidth]{importance_by_emotion.png}
\caption{Per-emotion concept importance. Red indicates that deleting the concept
reduces confidence (the concept supports the prediction); blue indicates that
deleting it increases confidence (the concept opposes the prediction).}
\label{fig:importance-emotion}
\end{figure}

The dominant concept, \texttt{long\_constant\_thick}, does not act uniformly. It
is the strongest positive contributor for \emph{sad} ($+1.05$),
\emph{disgusted} ($+1.09$) and \emph{happy} ($+0.49$), yet the strongest
negative contributor for \emph{fearful} ($-0.86$) and a negative contributor for
\emph{calm} ($-0.19$). Sustained harmonic energy therefore functions as
discriminative evidence in both directions: its presence pushes the model toward
the sustained-phonation emotions and away from \emph{fearful}, whose
realisations are characteristically less steady. An aggregate analysis averages
these opposing roles toward zero and loses the effect entirely.

Two further patterns are visible. For \emph{angry}, the falling-tone concepts
separate by steepness in opposite directions: \texttt{long\_dropping\_steep\_thick}
supports the prediction ($+0.33$) while \texttt{long\_dropping\_flat\_thick}
opposes it ($-0.39$), consistent with anger being marked by sharp rather than
gradual pitch falls. For \emph{fearful}, eleven of the twelve concepts have
negative or negligible effect, indicating that the class is identified largely
by the \emph{absence} of the modelled spectrogram structures rather than by the
presence of any of them.

\subsection{Discussion}
\label{sec:ablation-discussion}

The ablation study supports a qualified conclusion. The concept set is
meaningful: the subspace it spans is strongly class-relevant relative to a
random subspace of equal dimension, and leave-one-out deletion identifies a
clear and acoustically interpretable ordering, dominated by sustained harmonic
energy. The per-emotion decomposition further recovers signed, emotion-specific
roles that are consistent with the phonetic character of the classes.

The TCAV sensitivity scores themselves, however, do not survive the faithfulness
test: their ranking performs no better than chance under insertion and deletion.
We attribute this to the near-collinearity of the pairwise-fitted CAVs. This
suggests that concept-based explanation of this model should report ablation
importance rather than TCAV magnitude, and that the concept vocabulary would
benefit from being fitted so as to be mutually discriminative rather than merely
distinguishable from a random baseline.
 from the main document.
% Requires: booktabs, graphicx, xcolor (all standard).
% Bibliography entries are in concept_ablation.bib.

\section{Concept Importance via Insertion and Deletion}
\label{sec:concept-ablation}

TCAV~\cite{kim2018tcav} assigns each concept a sensitivity score, but a score is
not by itself evidence that the model \emph{uses} the concept. Insertion and
deletion~\cite{petsiuk2018rise} supplies that evidence: an explanation is
faithful if removing what it deems important degrades the prediction quickly,
and if restoring it to an uninformative baseline recovers the prediction
quickly. We adapt this protocol from input regions to concept directions.

\subsection{Method}
\label{sec:ablation-method}

Let $a \in \mathbb{R}^{D}$ be the activation of a sample at the layer in which
the Concept Activation Vectors (CAVs) were fitted, and let $c$ be the class the
model predicts for that sample. We orthonormalise the $K=12$ CAVs into a basis
$q_1,\dots,q_K$; orthonormalisation matters because the raw CAVs are strongly
correlated, and without it removing ``the top $k$ concepts'' would subtract a
shared component more than once.

For a ranking of the concepts, deletion and insertion at step $k$ are
\begin{align}
a^{\text{del}}_k &= a - \textstyle\sum_{i \in \mathcal{S}_k} (a \cdot q_i)\, q_i,\\
a^{\text{ins}}_k &= a - \textstyle\sum_{i=1}^{K} (a \cdot q_i)\, q_i
                     + \textstyle\sum_{i \in \mathcal{S}_k} (a \cdot q_i)\, q_i,
\end{align}
where $\mathcal{S}_k$ holds the $k$ highest-ranked concepts. By construction
$a^{\text{del}}_0 = a^{\text{ins}}_K = a$, so the two curves are complementary.
The modified activation is propagated through the remainder of the network and
we record the model's confidence in class $c$.

\paragraph{Confidence measure.}
We report the \emph{logit margin}, the logit of class $c$ minus the largest
competing logit, rather than the softmax probability. The classifier is
saturated on its own predictions (mean probability $0.99$), so probability-based
curves are flat to four decimal places regardless of the intervention; the
margin is a monotone transform of the same quantity but is not compressed.

\paragraph{Controls.}
Two controls separate two distinct claims. The \textsc{random} control applies
the identical procedure to the same concept subspace in shuffled order, testing
whether the \emph{ranking} carries information. The \textsc{noise} control
replaces the CAVs with $K$ random orthonormal directions, testing whether the
\emph{subspace} carries information or whether any $K$-dimensional subspace
would behave the same. All experiments use the $1{,}335$ correctly classified
RAVDESS recordings.

\subsection{Faithfulness of the TCAV ranking}
\label{sec:ablation-faithfulness}

Table~\ref{tab:ablation-auc} reports the area under the deletion and insertion
curves. A faithful ranking yields a low deletion AUC and a high insertion AUC.

\begin{table}[t]
\centering
\caption{Insertion--deletion areas under the curve, and the total confidence
lost when all $K=12$ directions are removed. Lower deletion AUC and higher
insertion AUC indicate a more faithful ranking.}
\label{tab:ablation-auc}
\begin{tabular}{lccc}
\toprule
Ordering & Deletion AUC $\downarrow$ & Insertion AUC $\uparrow$ & Total drop \\
\midrule
TCAV importance      & 8.82 & 8.67 & $-0.769$ \\
Shuffled (control)   & 8.75 & 8.75 & $-0.769$ \\
Random subspace (noise) & 9.09 & 9.09 & $-0.002$ \\
\bottomrule
\end{tabular}
\end{table}

Two conclusions follow, and they point in opposite directions.

The concepts are informative \emph{as a set}. Deleting the $12$ concept
directions costs $0.769$ of margin, whereas deleting $12$ random directions
costs $0.002$, a factor of roughly $380$. Since both interventions remove
exactly $12$ of $D=71{,}680$ dimensions, the difference is attributable
entirely to \emph{which} subspace is removed: the CAVs are aligned with
class-relevant structure rather than with arbitrary directions.

The TCAV ranking, however, is not faithful. Ordering deletions by TCAV
importance (AUC $8.82$) is no better than shuffling them (AUC $8.75$), and both
necessarily converge to the same endpoint because they remove the same subspace.
Were the per-sample importance scores meaningful, removing the highest-ranked
concepts first would degrade the prediction faster than removing them in
arbitrary order. It does not.

This is consistent with the geometry of the CAVs themselves: fitted pairwise
against a random counterexample set, the twelve directions have a mean pairwise
cosine similarity of $0.995$, so they span approximately one effective
dimension. A ranking cannot recover per-concept structure that the
representation does not contain.

\subsection{Magnitude of the effect}
\label{sec:ablation-percent}

The areas in Table~\ref{tab:ablation-auc} are in logit units, which are hard to
interpret on their own. Normalising every curve by the untouched decision margin
expresses the same result as a percentage of the model's confidence
(Figure~\ref{fig:percent-impact}).

Removing all twelve concept directions costs $8.5\%$ of the decision margin;
removing twelve random directions costs $0.02\%$. The concepts are therefore
roughly $380$ times more consequential than an arbitrary subspace of the same
dimension, while still accounting for under a tenth of the evidence the model
uses. This is the quantitative form of the qualified conclusion above: the
concept vocabulary is real but partial, and the remaining $91.5\%$ of the margin
rests on structure the twelve hand-designed patterns do not describe.

\begin{figure*}[t]
\centering
\includegraphics[width=\linewidth]{percent_impact.png}
\caption{Insertion and deletion expressed as a percentage of the original
decision margin. Left: deleting concept directions, most important first.
Centre: restoring them to a stripped activation. Right: the share of the margin
carried by the concepts, per emotion. Negative values indicate emotions whose
confidence \emph{increases} when the concepts are removed.}
\label{fig:percent-impact}
\end{figure*}

\begin{table}[t]
\centering
\caption{Share of the decision margin carried by the twelve concepts, per
emotion (ratio of means; $n$ is the number of correctly classified recordings).
A negative value means the concepts collectively argue \emph{against} the
predicted class.}
\label{tab:percent-emotion}
\begin{tabular}{lrrrr}
\toprule
Emotion & $n$ & Baseline & After removal & \% of margin \\
\midrule
sad       & 168 & 8.44 & 5.74 & $+32.0$ \\
disgusted & 182 & 9.09 & 6.72 & $+26.1$ \\
happy     & 169 & 9.14 & 8.54 & $+6.7$ \\
surprised & 180 & 8.67 & 8.09 & $+6.7$ \\
angry     & 182 & 8.57 & 8.01 & $+6.5$ \\
neutral   &  92 & 9.14 & 8.55 & $+6.4$ \\
calm      & 184 & 9.69 & 9.47 & $+2.3$ \\
fearful   & 178 & 9.94 & 11.41 & $-14.7$ \\
\midrule
overall   & 1335 & 9.09 & 8.32 & $+8.5$ \\
\bottomrule
\end{tabular}
\end{table}

The per-emotion split in Table~\ref{tab:percent-emotion} is markedly uneven. The
concepts carry a third of the margin for \emph{sad} and a quarter for
\emph{disgusted}, but only a few percent for \emph{calm}, \emph{angry},
\emph{neutral}, \emph{happy} and \emph{surprised}. For \emph{fearful} the effect
is negative: removing the concepts \emph{raises} the margin by $14.7\%$, so on
balance the modelled patterns are evidence against that class rather than for
it. The vocabulary is thus most explanatory for the sustained, low-arousal
classes and least explanatory---indeed counter-explanatory---for \emph{fearful}.

\subsection{Per-concept importance}
\label{sec:ablation-per-concept}

Because the ranking is uninformative, we attribute importance directly by
leave-one-out deletion: each concept is removed on its own and we record the
resulting change in margin. Table~\ref{tab:concept-importance} gives the result
aggregated over all emotions.

\begin{table}[t]
\centering
\caption{Leave-one-out concept importance, averaged over all correctly
classified recordings. A single random direction is included for scale.}
\label{tab:concept-importance}
\begin{tabular}{lcc}
\toprule
Concept & Margin drop & Share of total \\
\midrule
\texttt{long\_constant\_thick}        & 0.247 & 53.2\% \\
\texttt{long\_dropping\_steep\_thick} & 0.089 & 19.2\% \\
\texttt{short\_dropping\_steep\_thin} & 0.061 & 13.2\% \\
\texttt{short\_rising\_steep\_thick}  & 0.016 & 3.4\% \\
\texttt{long\_dropping\_flat\_thick}  & 0.014 & 3.0\% \\
\texttt{short\_rising\_steep\_thin}   & 0.010 & 2.1\% \\
\texttt{long\_dropping\_steep\_thin}  & 0.008 & 1.7\% \\
\texttt{short\_dropping\_steep\_thick}& 0.006 & 1.3\% \\
\texttt{long\_rising\_flat\_thick}    & 0.006 & 1.3\% \\
\texttt{long\_rising\_steep\_thin}    & 0.003 & 0.7\% \\
\texttt{long\_rising\_steep\_thick}   & 0.003 & 0.6\% \\
\texttt{short\_constant\_thick}       & 0.002 & 0.4\% \\
\midrule
\emph{single random direction}        & 0.0004 & --- \\
\bottomrule
\end{tabular}
\end{table}

Importance is heavily concentrated. \texttt{long\_constant\_thick}, a long
horizontal band corresponding to sustained harmonic energy, accounts for over
half of the total effect; the three highest-ranked concepts account for $85\%$.
The six lowest-ranked concepts contribute under $6\%$ combined, and the weakest
is within an order of magnitude of a random direction. We note that the
individual drops sum to $0.46$ while removing all twelve jointly costs $0.74$:
the concepts interact through the network rather than contributing additively,
so these values constitute a ranking rather than a decomposition.

\subsection{Per-emotion analysis}
\label{sec:ablation-per-emotion}

Aggregating over emotions conceals the more informative structure. Repeating the
leave-one-out analysis within each emotion (Table~\ref{tab:importance-emotion},
Figure~\ref{fig:importance-emotion}) shows that a concept's effect is
\emph{signed}: a positive drop means the concept supported the prediction, and a
negative drop means it argued against it, so that removing it increased the
model's confidence.

\begin{table}[t]
\centering
\small
\setlength{\tabcolsep}{3.5pt}
\caption{Change in logit margin when each concept is deleted, by emotion.
Positive values indicate the concept supports the prediction; negative values
indicate it opposes it. Concept names are abbreviated: L/S = long/short,
const/drop/rise = constant/dropping/rising, st/fl = steep/flat,
tk/tn = thick/thin.}
\label{tab:importance-emotion}
\begin{tabular}{lrrrrrrrrrrrr}
\toprule
& \rotatebox{90}{L-const/tk} & \rotatebox{90}{L-drop/fl/tk} & \rotatebox{90}{L-drop/st/tk}
& \rotatebox{90}{L-drop/st/tn} & \rotatebox{90}{L-rise/fl/tk} & \rotatebox{90}{L-rise/st/tk}
& \rotatebox{90}{L-rise/st/tn} & \rotatebox{90}{S-const/tk} & \rotatebox{90}{S-drop/st/tk}
& \rotatebox{90}{S-drop/st/tn} & \rotatebox{90}{S-rise/st/tk} & \rotatebox{90}{S-rise/st/tn} \\
\midrule
angry     &  0.27 & $-0.39$ &  0.33 &  0.04 & $-0.15$ & $-0.02$ &  0.02 & $-0.06$ & $-0.10$ &  0.25 &  0.09 & $-0.00$ \\
calm      & $-0.19$ &  0.05 & $-0.18$ &  0.02 &  0.20 & $-0.00$ & $-0.00$ & $-0.02$ &  0.08 & $-0.06$ &  0.01 &  0.04 \\
disgusted &  1.09 &  0.33 &  0.23 &  0.00 &  0.15 &  0.05 &  0.00 &  0.07 &  0.06 &  0.11 & $-0.05$ & $-0.02$ \\
fearful   & $-0.86$ & $-0.37$ & $-0.26$ & $-0.00$ & $-0.15$ & $-0.04$ & $-0.00$ & $-0.07$ & $-0.13$ &  0.14 &  0.06 &  0.02 \\
happy     &  0.49 & $-0.19$ &  0.01 & $-0.01$ & $-0.12$ & $-0.02$ &  0.01 &  0.01 & $-0.03$ &  0.05 &  0.04 &  0.01 \\
neutral   &  0.07 &  0.08 &  0.10 & $-0.00$ & $-0.11$ & $-0.03$ &  0.00 &  0.03 &  0.15 &  0.10 & $-0.04$ &  0.02 \\
sad       &  1.05 &  0.65 &  0.24 &  0.03 &  0.34 &  0.07 & $-0.00$ &  0.05 &  0.09 & $-0.22$ &  0.02 &  0.02 \\
surprised &  0.15 &  0.08 &  0.18 & $-0.00$ & $-0.01$ &  0.02 & $-0.00$ &  0.01 & $-0.03$ &  0.05 & $-0.02$ & $-0.02$ \\
\bottomrule
\end{tabular}
\end{table}

\begin{figure}[t]
\centering
\includegraphics[width=\linewidth]{importance_by_emotion.png}
\caption{Per-emotion concept importance. Red indicates that deleting the concept
reduces confidence (the concept supports the prediction); blue indicates that
deleting it increases confidence (the concept opposes the prediction).}
\label{fig:importance-emotion}
\end{figure}

The dominant concept, \texttt{long\_constant\_thick}, does not act uniformly. It
is the strongest positive contributor for \emph{sad} ($+1.05$),
\emph{disgusted} ($+1.09$) and \emph{happy} ($+0.49$), yet the strongest
negative contributor for \emph{fearful} ($-0.86$) and a negative contributor for
\emph{calm} ($-0.19$). Sustained harmonic energy therefore functions as
discriminative evidence in both directions: its presence pushes the model toward
the sustained-phonation emotions and away from \emph{fearful}, whose
realisations are characteristically less steady. An aggregate analysis averages
these opposing roles toward zero and loses the effect entirely.

Two further patterns are visible. For \emph{angry}, the falling-tone concepts
separate by steepness in opposite directions: \texttt{long\_dropping\_steep\_thick}
supports the prediction ($+0.33$) while \texttt{long\_dropping\_flat\_thick}
opposes it ($-0.39$), consistent with anger being marked by sharp rather than
gradual pitch falls. For \emph{fearful}, eleven of the twelve concepts have
negative or negligible effect, indicating that the class is identified largely
by the \emph{absence} of the modelled spectrogram structures rather than by the
presence of any of them.

\subsection{Discussion}
\label{sec:ablation-discussion}

The ablation study supports a qualified conclusion. The concept set is
meaningful: the subspace it spans is strongly class-relevant relative to a
random subspace of equal dimension, and leave-one-out deletion identifies a
clear and acoustically interpretable ordering, dominated by sustained harmonic
energy. The per-emotion decomposition further recovers signed, emotion-specific
roles that are consistent with the phonetic character of the classes.

The TCAV sensitivity scores themselves, however, do not survive the faithfulness
test: their ranking performs no better than chance under insertion and deletion.
We attribute this to the near-collinearity of the pairwise-fitted CAVs. This
suggests that concept-based explanation of this model should report ablation
importance rather than TCAV magnitude, and that the concept vocabulary would
benefit from being fitted so as to be mutually discriminative rather than merely
distinguishable from a random baseline.
 from the main document.
% Requires: booktabs, graphicx, xcolor (all standard).
% Bibliography entries are in concept_ablation.bib.

\section{Concept Importance via Insertion and Deletion}
\label{sec:concept-ablation}

TCAV~\cite{kim2018tcav} assigns each concept a sensitivity score, but a score is
not by itself evidence that the model \emph{uses} the concept. Insertion and
deletion~\cite{petsiuk2018rise} supplies that evidence: an explanation is
faithful if removing what it deems important degrades the prediction quickly,
and if restoring it to an uninformative baseline recovers the prediction
quickly. We adapt this protocol from input regions to concept directions.

\subsection{Method}
\label{sec:ablation-method}

Let $a \in \mathbb{R}^{D}$ be the activation of a sample at the layer in which
the Concept Activation Vectors (CAVs) were fitted, and let $c$ be the class the
model predicts for that sample. We orthonormalise the $K=12$ CAVs into a basis
$q_1,\dots,q_K$; orthonormalisation matters because the raw CAVs are strongly
correlated, and without it removing ``the top $k$ concepts'' would subtract a
shared component more than once.

For a ranking of the concepts, deletion and insertion at step $k$ are
\begin{align}
a^{\text{del}}_k &= a - \textstyle\sum_{i \in \mathcal{S}_k} (a \cdot q_i)\, q_i,\\
a^{\text{ins}}_k &= a - \textstyle\sum_{i=1}^{K} (a \cdot q_i)\, q_i
                     + \textstyle\sum_{i \in \mathcal{S}_k} (a \cdot q_i)\, q_i,
\end{align}
where $\mathcal{S}_k$ holds the $k$ highest-ranked concepts. By construction
$a^{\text{del}}_0 = a^{\text{ins}}_K = a$, so the two curves are complementary.
The modified activation is propagated through the remainder of the network and
we record the model's confidence in class $c$.

\paragraph{Confidence measure.}
We report the \emph{logit margin}, the logit of class $c$ minus the largest
competing logit, rather than the softmax probability. The classifier is
saturated on its own predictions (mean probability $0.99$), so probability-based
curves are flat to four decimal places regardless of the intervention; the
margin is a monotone transform of the same quantity but is not compressed.

\paragraph{Controls.}
Two controls separate two distinct claims. The \textsc{random} control applies
the identical procedure to the same concept subspace in shuffled order, testing
whether the \emph{ranking} carries information. The \textsc{noise} control
replaces the CAVs with $K$ random orthonormal directions, testing whether the
\emph{subspace} carries information or whether any $K$-dimensional subspace
would behave the same. All experiments use the $1{,}335$ correctly classified
RAVDESS recordings.

\subsection{Faithfulness of the TCAV ranking}
\label{sec:ablation-faithfulness}

Table~\ref{tab:ablation-auc} reports the area under the deletion and insertion
curves. A faithful ranking yields a low deletion AUC and a high insertion AUC.

\begin{table}[t]
\centering
\caption{Insertion--deletion areas under the curve, and the total confidence
lost when all $K=12$ directions are removed. Lower deletion AUC and higher
insertion AUC indicate a more faithful ranking.}
\label{tab:ablation-auc}
\begin{tabular}{lccc}
\toprule
Ordering & Deletion AUC $\downarrow$ & Insertion AUC $\uparrow$ & Total drop \\
\midrule
TCAV importance      & 8.82 & 8.67 & $-0.769$ \\
Shuffled (control)   & 8.75 & 8.75 & $-0.769$ \\
Random subspace (noise) & 9.09 & 9.09 & $-0.002$ \\
\bottomrule
\end{tabular}
\end{table}

Two conclusions follow, and they point in opposite directions.

The concepts are informative \emph{as a set}. Deleting the $12$ concept
directions costs $0.769$ of margin, whereas deleting $12$ random directions
costs $0.002$, a factor of roughly $380$. Since both interventions remove
exactly $12$ of $D=71{,}680$ dimensions, the difference is attributable
entirely to \emph{which} subspace is removed: the CAVs are aligned with
class-relevant structure rather than with arbitrary directions.

The TCAV ranking, however, is not faithful. Ordering deletions by TCAV
importance (AUC $8.82$) is no better than shuffling them (AUC $8.75$), and both
necessarily converge to the same endpoint because they remove the same subspace.
Were the per-sample importance scores meaningful, removing the highest-ranked
concepts first would degrade the prediction faster than removing them in
arbitrary order. It does not.

This is consistent with the geometry of the CAVs themselves: fitted pairwise
against a random counterexample set, the twelve directions have a mean pairwise
cosine similarity of $0.995$, so they span approximately one effective
dimension. A ranking cannot recover per-concept structure that the
representation does not contain.

\subsection{Magnitude of the effect}
\label{sec:ablation-percent}

The areas in Table~\ref{tab:ablation-auc} are in logit units, which are hard to
interpret on their own. Normalising every curve by the untouched decision margin
expresses the same result as a percentage of the model's confidence
(Figure~\ref{fig:percent-impact}).

Removing all twelve concept directions costs $8.5\%$ of the decision margin;
removing twelve random directions costs $0.02\%$. The concepts are therefore
roughly $380$ times more consequential than an arbitrary subspace of the same
dimension, while still accounting for under a tenth of the evidence the model
uses. This is the quantitative form of the qualified conclusion above: the
concept vocabulary is real but partial, and the remaining $91.5\%$ of the margin
rests on structure the twelve hand-designed patterns do not describe.

\begin{figure*}[t]
\centering
\includegraphics[width=\linewidth]{percent_impact.png}
\caption{Insertion and deletion expressed as a percentage of the original
decision margin. Left: deleting concept directions, most important first.
Centre: restoring them to a stripped activation. Right: the share of the margin
carried by the concepts, per emotion. Negative values indicate emotions whose
confidence \emph{increases} when the concepts are removed.}
\label{fig:percent-impact}
\end{figure*}

\begin{table}[t]
\centering
\caption{Share of the decision margin carried by the twelve concepts, per
emotion (ratio of means; $n$ is the number of correctly classified recordings).
A negative value means the concepts collectively argue \emph{against} the
predicted class.}
\label{tab:percent-emotion}
\begin{tabular}{lrrrr}
\toprule
Emotion & $n$ & Baseline & After removal & \% of margin \\
\midrule
sad       & 168 & 8.44 & 5.74 & $+32.0$ \\
disgusted & 182 & 9.09 & 6.72 & $+26.1$ \\
happy     & 169 & 9.14 & 8.54 & $+6.7$ \\
surprised & 180 & 8.67 & 8.09 & $+6.7$ \\
angry     & 182 & 8.57 & 8.01 & $+6.5$ \\
neutral   &  92 & 9.14 & 8.55 & $+6.4$ \\
calm      & 184 & 9.69 & 9.47 & $+2.3$ \\
fearful   & 178 & 9.94 & 11.41 & $-14.7$ \\
\midrule
overall   & 1335 & 9.09 & 8.32 & $+8.5$ \\
\bottomrule
\end{tabular}
\end{table}

The per-emotion split in Table~\ref{tab:percent-emotion} is markedly uneven. The
concepts carry a third of the margin for \emph{sad} and a quarter for
\emph{disgusted}, but only a few percent for \emph{calm}, \emph{angry},
\emph{neutral}, \emph{happy} and \emph{surprised}. For \emph{fearful} the effect
is negative: removing the concepts \emph{raises} the margin by $14.7\%$, so on
balance the modelled patterns are evidence against that class rather than for
it. The vocabulary is thus most explanatory for the sustained, low-arousal
classes and least explanatory---indeed counter-explanatory---for \emph{fearful}.

\subsection{Per-concept importance}
\label{sec:ablation-per-concept}

Because the ranking is uninformative, we attribute importance directly by
leave-one-out deletion: each concept is removed on its own and we record the
resulting change in margin. Table~\ref{tab:concept-importance} gives the result
aggregated over all emotions.

\begin{table}[t]
\centering
\caption{Leave-one-out concept importance, averaged over all correctly
classified recordings. A single random direction is included for scale.}
\label{tab:concept-importance}
\begin{tabular}{lcc}
\toprule
Concept & Margin drop & Share of total \\
\midrule
\texttt{long\_constant\_thick}        & 0.247 & 53.2\% \\
\texttt{long\_dropping\_steep\_thick} & 0.089 & 19.2\% \\
\texttt{short\_dropping\_steep\_thin} & 0.061 & 13.2\% \\
\texttt{short\_rising\_steep\_thick}  & 0.016 & 3.4\% \\
\texttt{long\_dropping\_flat\_thick}  & 0.014 & 3.0\% \\
\texttt{short\_rising\_steep\_thin}   & 0.010 & 2.1\% \\
\texttt{long\_dropping\_steep\_thin}  & 0.008 & 1.7\% \\
\texttt{short\_dropping\_steep\_thick}& 0.006 & 1.3\% \\
\texttt{long\_rising\_flat\_thick}    & 0.006 & 1.3\% \\
\texttt{long\_rising\_steep\_thin}    & 0.003 & 0.7\% \\
\texttt{long\_rising\_steep\_thick}   & 0.003 & 0.6\% \\
\texttt{short\_constant\_thick}       & 0.002 & 0.4\% \\
\midrule
\emph{single random direction}        & 0.0004 & --- \\
\bottomrule
\end{tabular}
\end{table}

Importance is heavily concentrated. \texttt{long\_constant\_thick}, a long
horizontal band corresponding to sustained harmonic energy, accounts for over
half of the total effect; the three highest-ranked concepts account for $85\%$.
The six lowest-ranked concepts contribute under $6\%$ combined, and the weakest
is within an order of magnitude of a random direction. We note that the
individual drops sum to $0.46$ while removing all twelve jointly costs $0.74$:
the concepts interact through the network rather than contributing additively,
so these values constitute a ranking rather than a decomposition.

\subsection{Per-emotion analysis}
\label{sec:ablation-per-emotion}

Aggregating over emotions conceals the more informative structure. Repeating the
leave-one-out analysis within each emotion (Table~\ref{tab:importance-emotion},
Figure~\ref{fig:importance-emotion}) shows that a concept's effect is
\emph{signed}: a positive drop means the concept supported the prediction, and a
negative drop means it argued against it, so that removing it increased the
model's confidence.

\begin{table}[t]
\centering
\small
\setlength{\tabcolsep}{3.5pt}
\caption{Change in logit margin when each concept is deleted, by emotion.
Positive values indicate the concept supports the prediction; negative values
indicate it opposes it. Concept names are abbreviated: L/S = long/short,
const/drop/rise = constant/dropping/rising, st/fl = steep/flat,
tk/tn = thick/thin.}
\label{tab:importance-emotion}
\begin{tabular}{lrrrrrrrrrrrr}
\toprule
& \rotatebox{90}{L-const/tk} & \rotatebox{90}{L-drop/fl/tk} & \rotatebox{90}{L-drop/st/tk}
& \rotatebox{90}{L-drop/st/tn} & \rotatebox{90}{L-rise/fl/tk} & \rotatebox{90}{L-rise/st/tk}
& \rotatebox{90}{L-rise/st/tn} & \rotatebox{90}{S-const/tk} & \rotatebox{90}{S-drop/st/tk}
& \rotatebox{90}{S-drop/st/tn} & \rotatebox{90}{S-rise/st/tk} & \rotatebox{90}{S-rise/st/tn} \\
\midrule
angry     &  0.27 & $-0.39$ &  0.33 &  0.04 & $-0.15$ & $-0.02$ &  0.02 & $-0.06$ & $-0.10$ &  0.25 &  0.09 & $-0.00$ \\
calm      & $-0.19$ &  0.05 & $-0.18$ &  0.02 &  0.20 & $-0.00$ & $-0.00$ & $-0.02$ &  0.08 & $-0.06$ &  0.01 &  0.04 \\
disgusted &  1.09 &  0.33 &  0.23 &  0.00 &  0.15 &  0.05 &  0.00 &  0.07 &  0.06 &  0.11 & $-0.05$ & $-0.02$ \\
fearful   & $-0.86$ & $-0.37$ & $-0.26$ & $-0.00$ & $-0.15$ & $-0.04$ & $-0.00$ & $-0.07$ & $-0.13$ &  0.14 &  0.06 &  0.02 \\
happy     &  0.49 & $-0.19$ &  0.01 & $-0.01$ & $-0.12$ & $-0.02$ &  0.01 &  0.01 & $-0.03$ &  0.05 &  0.04 &  0.01 \\
neutral   &  0.07 &  0.08 &  0.10 & $-0.00$ & $-0.11$ & $-0.03$ &  0.00 &  0.03 &  0.15 &  0.10 & $-0.04$ &  0.02 \\
sad       &  1.05 &  0.65 &  0.24 &  0.03 &  0.34 &  0.07 & $-0.00$ &  0.05 &  0.09 & $-0.22$ &  0.02 &  0.02 \\
surprised &  0.15 &  0.08 &  0.18 & $-0.00$ & $-0.01$ &  0.02 & $-0.00$ &  0.01 & $-0.03$ &  0.05 & $-0.02$ & $-0.02$ \\
\bottomrule
\end{tabular}
\end{table}

\begin{figure}[t]
\centering
\includegraphics[width=\linewidth]{importance_by_emotion.png}
\caption{Per-emotion concept importance. Red indicates that deleting the concept
reduces confidence (the concept supports the prediction); blue indicates that
deleting it increases confidence (the concept opposes the prediction).}
\label{fig:importance-emotion}
\end{figure}

The dominant concept, \texttt{long\_constant\_thick}, does not act uniformly. It
is the strongest positive contributor for \emph{sad} ($+1.05$),
\emph{disgusted} ($+1.09$) and \emph{happy} ($+0.49$), yet the strongest
negative contributor for \emph{fearful} ($-0.86$) and a negative contributor for
\emph{calm} ($-0.19$). Sustained harmonic energy therefore functions as
discriminative evidence in both directions: its presence pushes the model toward
the sustained-phonation emotions and away from \emph{fearful}, whose
realisations are characteristically less steady. An aggregate analysis averages
these opposing roles toward zero and loses the effect entirely.

Two further patterns are visible. For \emph{angry}, the falling-tone concepts
separate by steepness in opposite directions: \texttt{long\_dropping\_steep\_thick}
supports the prediction ($+0.33$) while \texttt{long\_dropping\_flat\_thick}
opposes it ($-0.39$), consistent with anger being marked by sharp rather than
gradual pitch falls. For \emph{fearful}, eleven of the twelve concepts have
negative or negligible effect, indicating that the class is identified largely
by the \emph{absence} of the modelled spectrogram structures rather than by the
presence of any of them.

\subsection{Discussion}
\label{sec:ablation-discussion}

The ablation study supports a qualified conclusion. The concept set is
meaningful: the subspace it spans is strongly class-relevant relative to a
random subspace of equal dimension, and leave-one-out deletion identifies a
clear and acoustically interpretable ordering, dominated by sustained harmonic
energy. The per-emotion decomposition further recovers signed, emotion-specific
roles that are consistent with the phonetic character of the classes.

The TCAV sensitivity scores themselves, however, do not survive the faithfulness
test: their ranking performs no better than chance under insertion and deletion.
We attribute this to the near-collinearity of the pairwise-fitted CAVs. This
suggests that concept-based explanation of this model should report ablation
importance rather than TCAV magnitude, and that the concept vocabulary would
benefit from being fitted so as to be mutually discriminative rather than merely
distinguishable from a random baseline.
